\documentclass[12pt,a4paper]{article}
\usepackage[utf8]{inputenc}
\usepackage{amsmath, amssymb}
\usepackage{geometry}
\usepackage{setspace}
\usepackage{graphicx}

\geometry{margin=2.5cm}
\setstretch{1.2}

\title{Taller de Regla de L'Hopital}
\author{Jose Luciano Mejia Arias}
\date{}

\makeatletter
\renewcommand{\maketitle}{
  \begin{flushright}
    \includegraphics[width=3cm]{../../logo_unal.png}
  \end{flushright}
  \begin{center}
    {\LARGE \textbf{\@title}}\\[0.3cm]
    {\large \@author}
  \end{center}
}
\makeatother

\begin{document}
\maketitle

\section*{Instrucciones}
Resuelva cada ejercicio mostrando el desarrollo. Indique cuando use la regla de L'Hopital y verifique que se cumplan sus condiciones antes de aplicarla.

\section{Límites con forma $0/0$}
Calcule los siguientes límites:

\begin{enumerate}
    \item $\displaystyle \lim_{x\to 0} \frac{\sin(4x)}{x}$
    \item $\displaystyle \lim_{x\to 0} \frac{e^{2x}-1}{x}$
    \item $\displaystyle \lim_{x\to 0} \frac{1-\cos(2x)}{x^2}$
    \item $\displaystyle \lim_{x\to 0} \frac{\sqrt{1+x}-1}{x}$
    \item $\displaystyle \lim_{x\to 2} \frac{x^2-4}{x-2}$
\end{enumerate}

\section{Límites con forma $\infty/\infty$}
Calcule los siguientes límites:

\begin{enumerate}
    \item $\displaystyle \lim_{x\to \infty} \frac{\ln x}{x}$
    \item $\displaystyle \lim_{x\to \infty} \frac{x}{e^x}$
    \item $\displaystyle \lim_{x\to \infty} \frac{x^2}{e^x}$
    \item $\displaystyle \lim_{x\to \infty} \frac{5x-1}{2x+7}$
\end{enumerate}

\section{Aplicaciones con varias derivadas}
Resuelva los siguientes límites, aplicando L'Hopital las veces que sea necesario:

\begin{enumerate}
    \item $\displaystyle \lim_{x\to 0} \frac{x-\sin x}{x^3}$
    \item $\displaystyle \lim_{x\to 0} \frac{\tan x - x}{x^3}$
    \item $\displaystyle \lim_{x\to 0} \frac{e^x - 1 - x}{x^2}$
    \item $\displaystyle \lim_{x\to 0} \frac{\ln(1+x)}{x}$
    \item $\displaystyle \lim_{x\to \infty} \frac{\ln x}{\sqrt{x}}$
\end{enumerate}

\section{Ejercicios de aplicación}

\begin{enumerate}
    \item Halle la pendiente de la recta tangente a $f(x)=\dfrac{\sin x}{x}$ en $x=0$ usando el límite correspondiente.
    \item Determine si la función $g(x)=\dfrac{e^x}{x^2}$ tiende a $0$, $\infty$ o a un valor finito cuando $x\to\infty$.
    \item Use L'Hopital para justificar que $\displaystyle \lim_{x\to 0} \frac{\ln(1+x)}{x}=1$.
\end{enumerate}

\end{document}